% Part: turing-machines
% Chapter: machines-computations
% Section: unsolvability-decision-problem

\documentclass[../../../include/open-logic-section]{subfiles}

\begin{document}

\olfileid{tur}{und}{dec}
\olsection{সিদ্ধান্ত সমস্যা}

প্রদত্ত কোনো !!{sentence} সিদ্ধ কি না নির্ধারণ করার কার্যকর পদ্ধতি থাকলে
এবং কেবল তখনই প্রথম-ক্রমের যুক্তিবিদ্যাকে \emph{নির্ণেয়} বলি। বাস্তবে এমন
কোনো পদ্ধতি নেই: প্রথম-ক্রমের বাক্যের সিদ্ধতা নির্ণয় করার সমস্যাটি
অসমাধানযোগ্য।

এই গুরুত্বপূর্ণ নেতিবাচক ফল প্রতিষ্ঠা করতে আমরা প্রমাণ করব যে সিদ্ধান্ত
সমস্যাটি কোনো টুরিং যন্ত্রে সমাধান করা যায় না। অর্থাৎ আমরা দেখাব, ফিতায়
প্রথম-ক্রমের কোনো !!{sentence} নিয়ে শুরু করলে ওই !!{sentence} সিদ্ধ কি না অনুযায়ী
শেষ পর্যন্ত থেমে নির্গম হিসেবে $1$ অথবা~$0$ দেয়—এমন কোনো টুরিং যন্ত্র
নেই। চার্চ--টুরিং থিসিস অনুসারে প্রতিটি গণনসাধ্য অপেক্ষক টুরিং-গণনসাধ্য।
তাই এই “সিদ্ধতা-অপেক্ষক” আদৌ কার্যকরভাবে গণনসাধ্য হলে টুরিং-গণনসাধ্যও
হতো। এটি টুরিং-গণনসাধ্য না হলে কার্যকরভাবেও গণনসাধ্য হতে পারে না।

সিদ্ধান্ত সমস্যার অসমাধানযোগ্যতা প্রমাণের কৌশল হলো থামার সমস্যাকে তাতে
হ্রাস করা। এর অর্থ নিচের কথা: আমরা প্রমাণ করেছি, সূচক~$e$-এর টুরিং
যন্ত্রটি নিবেশ~$w$-তে থামলে যে অপেক্ষক~$h(e,w)$-এর মান~$1$, আর
অন্যথায় যার মান~$0$, সেটি টুরিং-গণনসাধ্য নয়। আমরা দেখাব, প্রথম-ক্রমের
বাক্যের সিদ্ধতা নির্ণয়কারী কোনো টুরিং যন্ত্র থাকলে~$h$ গণনাকারী টুরিং
যন্ত্রও থাকত। যেহেতু কোনো টুরিং যন্ত্র~$h$ গণনা করতে পারে না, তাই সিদ্ধতা
নির্ণয়কারী টুরিং যন্ত্রও থাকতে পারে না।

এই কৌশলের প্রথম ধাপে দেখাতে হবে, প্রতিটি নিবেশ~$w$ ও টুরিং যন্ত্র~$M$-এর
জন্য যন্ত্র~$M$-এর নির্দেশ-সেট ও নিবেশ~$w$ প্রকাশকারী !!a{sentence}
$!T(M, w)$ এবং “$M$ শেষ পর্যন্ত থামে” প্রকাশকারী !!a{sentence}~$!E(M, w)$
কার্যকরভাবে বর্ণনা করা যায়, যেন:
\begin{quote}
  $\Entails !T(M, w) \lif !E(M,w)$ যদি এবং কেবল যদি $M$ নিবেশ~$w$-তে থামে।
\end{quote}
আমাদের প্রমাণের প্রধান অংশে এই বাক্যদুটি~$!T(M, w)$ ও~$!E(M, w)$ বর্ণনা
করা হবে এবং যাচাই করা হবে যে~$M$ নিবেশ~$w$-তে থামলে এবং কেবল তখনই
$!T(M, w) \lif !E(M, w)$ সিদ্ধ।

\end{document}
